
\documentclass[conference]{IEEEtran}
\IEEEoverridecommandlockouts

\usepackage{cite}
\usepackage{amsmath,amssymb}
\usepackage{graphicx}
\usepackage{booktabs}
\usepackage{hyperref}
\usepackage{xcolor}

\hypersetup{
    colorlinks=true,
    urlcolor=blue,
    citecolor=blue
}

\title{Predicting 30-Day Hospital Readmissions
Using Explainable Machine Learning:
A Comparative Study}

\author{
\IEEEauthorblockN{
Bahaa Al-Juboori\IEEEauthorrefmark{1},
Shaaz E. Jahan\IEEEauthorrefmark{1},
Jason Fakidis\IEEEauthorrefmark{1}
}
\IEEEauthorblockA{
\IEEEauthorrefmark{1}
British Columbia Institute of Technology\\
Burnaby, BC, Canada
}
}

\begin{document}
\maketitle

\begin{abstract}
Hospital readmissions create significant challenges
for healthcare systems. Identifying patients at risk
of early readmission may support better discharge
planning and follow-up care.

This project proposes a comparative study of machine
learning methods for predicting hospital readmission
within 30 days using the Diabetes 130-US Hospitals
dataset. We plan to compare Logistic Regression,
Random Forest, XGBoost, and a Multi-Layer Perceptron
(MLP) neural network.

We will investigate the effects of data preprocessing,
class imbalance, and model selection on predictive
performance. Explainability techniques such as SHAP
will be considered to identify influential predictors.
The goal is to evaluate the trade-off between
predictive performance and interpretability rather
than assuming that more complex models are superior.
\end{abstract}

\begin{IEEEkeywords}
Hospital Readmission, Machine Learning,
Healthcare Analytics, XGBoost, Neural Networks,
Explainable AI, SHAP
\end{IEEEkeywords}

\section{Problem Definition and Motivation}

Unplanned hospital readmissions can place additional
pressure on healthcare resources and may indicate
the need for improved post-discharge support.

Our primary research question is:

\textit{How do traditional machine learning methods
and neural networks compare in predicting 30-day
hospital readmissions among patients with diabetes?}

We will also investigate whether explainability
methods provide useful insights into model
predictions.

\section{Preliminary Literature Review}

Strack et al.~\cite{strack2014} investigated factors
associated with hospital readmission among patients
with diabetes using a large collection of hospital
encounters.

Their work provides the basis for the dataset
selected for this project.

Class imbalance presents an additional challenge,
because positive readmission cases may represent
a minority of observations. Chawla et
al.~\cite{chawla2002} introduced SMOTE, a widely
used technique for addressing imbalanced
classification datasets.

Lundberg and Lee~\cite{lundberg2017} introduced
SHAP, an approach for interpreting model
predictions by estimating feature contributions.

These works motivate our comparison of predictive
models, imbalance-handling methods, and
interpretability techniques.

\section{Dataset}

We plan to use the publicly available
Diabetes 130-US Hospitals dataset from the
UCI Machine Learning Repository.

The dataset contains approximately 101,766
hospital encounters collected between 1999
and 2008.

Potential predictors include:
\begin{itemize}
    \item Patient age group
    \item Length of hospital stay
    \item Previous inpatient and emergency visits
    \item Number of diagnoses
    \item Medication-related variables
\end{itemize}

We will define the binary target as:

\begin{equation}
y =
\begin{cases}
1, & \text{readmission within 30 days}\\
0, & \text{otherwise}
\end{cases}
\end{equation}

We will examine missing values, duplicate
encounters, feature availability, and class
distributions before training models.

\section{Methods Under Study}

We propose four supervised classification models:

\begin{enumerate}
    \item Logistic Regression as an interpretable
    baseline.
    \item Random Forest to capture nonlinear
    relationships.
    \item XGBoost for gradient-boosted decision trees.
    \item Multi-Layer Perceptron (MLP) as a
    neural-network comparison.
\end{enumerate}

All models will be evaluated using a consistent
data partition and preprocessing protocol.

Preprocessing will include appropriate missing-value
treatment, categorical encoding, and feature scaling
where required.

To reduce leakage, preprocessing will be fitted
only using training data. We will also investigate
whether repeated encounters from the same patient
require patient-level data splitting.

\section{Evaluation Methodology}

Because early readmissions may be relatively
uncommon, classification accuracy alone will not
be considered sufficient.

We plan to report:
\begin{itemize}
    \item Precision and Recall
    \item F1-score
    \item ROC-AUC
    \item Precision-Recall AUC
    \item Confusion matrices
\end{itemize}

Class-weighting and resampling methods may be
evaluated using training data only.

For selected trained models, SHAP-based analysis
will be considered to explore feature contributions.

\section{Plan and Methodology}

Our proposed development phases are:

\begin{enumerate}
    \item Data acquisition and exploratory analysis
    \item Data cleaning and feature engineering
    \item Baseline model development
    \item Advanced machine learning experiments
    \item Evaluation and error analysis
    \item Model interpretation and final reporting
\end{enumerate}

All experiments will be documented with reproducible
Python notebooks and version-controlled source code.

\section{Team Organization}

The proposed division of responsibilities is:

\textbf{Bahaa Al-Juboori:}


\textbf{Shaaz E. Jahan:}


\textbf{Jason Fakidis:}


\section{Reproducibility and Repository}

Project materials will be maintained in a public
GitHub repository, including documentation,
Python notebooks, source code, and results.

\noindent\url{
https://github.com/bahaali273/COMP9170-Hospital-Readmission
}

Raw data will be obtained from the original
UCI source rather than redistributed through
the repository.

\section{Expected Outcomes}

We expect to produce a reproducible comparison
of classification methods for early hospital
readmission prediction.

The final study will discuss predictive performance,
class imbalance, model interpretability,
methodological limitations, and potential risks
associated with using historical clinical data.

No experimental performance claims are made
at the proposal stage.

\begin{thebibliography}{00}

\bibitem{strack2014}
B. Strack et al.,
``Impact of HbA1c Measurement on Hospital
Readmission Rates: Analysis of 70,000
Clinical Database Patient Records,''
\textit{BioMed Research International},
2014.

\bibitem{chawla2002}
N. V. Chawla et al.,
``SMOTE: Synthetic Minority
Over-sampling Technique,''
\textit{Journal of Artificial
Intelligence Research},
vol. 16, pp. 321--357, 2002.

\bibitem{lundberg2017}
S. M. Lundberg and S.-I. Lee,
``A Unified Approach to Interpreting
Model Predictions,''
in \textit{Advances in Neural Information
Processing Systems}, 2017.

\bibitem{uci}
UCI Machine Learning Repository,
``Diabetes 130-US Hospitals for Years
1999--2008.''
[Online]. Available:
\url{https://archive.ics.uci.edu/dataset/296/}

\end{thebibliography}

\end{document}
